\documentclass[10pt,a4paper]{amsart}
\usepackage[margin=19mm]{geometry}
\usepackage{amsmath,amssymb,mathtools}
\usepackage[scr=rsfs,cal=euler]{mathalfa}
\usepackage{xfrac}
\usepackage[unicode,hidelinks,bookmarks=false]{hyperref}
\newcommand{\ie}{i.e.\ }
\newcommand{\cf}{cf.\ }
\newcommand{\resp}{respectively\ }
\newcommand{\Subsection}{\subsection}
\providecommand{\nobreakdash}{\nobreak\hbox{-}}
\setlength{\emergencystretch}{2em}
\multlinegap0pt
\newtheorem{theorem}{Theorem}[section]
\newtheorem{lemma}[theorem]{Lemma}
\newtheorem{proposition}[theorem]{Proposition}
\newtheorem{definition}[theorem]{Definition}
\newtheorem{remark}[theorem]{Remark}
\newtheorem{example}[theorem]{Example}
\begin{document}
\title[Weyl Pseudo Almost Periodic Type Solutions to Semilinear Sto]{Weyl Pseudo Almost Periodic Type Solutions to Semilinear Stochastic Evolution Equations Driven by Fractional Brownian Motion}
\author{Dimplekumar N. Chalishajar, Marko Kostic, Daniel Velinov}
\date{}
\maketitle
\noindent\emph{Source and licence.} Weyl Pseudo Almost Periodic Type Solutions to Semilinear Stochastic Evolution Equations Driven by Fractional Brownian Motion. Version arXiv:2609.25321v1 (CC BY 4.0). This material is distributed under the Creative Commons Attribution 4.0 International licence, http://creativecommons.org/licenses/by/4.0/, as recorded in the arXiv OAI licence metadata for this exact version and shown on the abstract page https://arxiv.org/abs/2609.25321v1. Full rights notice: licenses/arxiv-2609.25321-CC-BY-4.0.txt.\par\smallskip
\noindent\textit{Editorial scope.} Counted unit: the original \texttt{theorem} environment \texttt{E2char} at source lines 232--240 together with its complete proof at lines 242--268; the delivered document keeps source lines 64--269 so that every referenced label and every citation resolves inside it. Native editable source is the paper's own concatenated TeX; the AMS wrapper is editorial formatting only. No publisher class, upstream script or unrelated artwork is redistributed.
\section{Preliminaries}

Throughout the paper, $(X,\|\cdot\|)$ denotes a separable real Hilbert space and $(\Omega,\mathcal F,p,\{\mathcal F_t\}_{t\in\mathbb R})$ a complete filtered probability space, where, following the usual convention, $p$ denotes the probability measure. We denote by $H:=L^2(\Omega,X)$ the space of $X$-valued random variables $x$ with $E\|x\|^2<\infty$, equipped with the inner product $\langle x,y\rangle_H:=E\langle x,y\rangle_X$ and norm $\|x\|_H:=(E\|x\|^2)^{1/2}$; $(H,\langle\cdot,\cdot\rangle_H)$ is itself a separable Hilbert space. We denote by $L^2_{loc}(\mathbb R,X)$, respectively $L^2_{loc}(\mathbb R,H)$, the corresponding spaces of locally square-integrable functions.

\subsection{The deterministic Weyl seminorm}

For $\varphi\in L^p_{loc}(\mathbb R,X)$, $1\le p<\infty$, the Weyl seminorm of $\varphi$ is defined by
\begin{align*}
\|\varphi\|_{W^p}=\lim_{\flat\to\infty}\sup_{\vartheta\in\mathbb R}\left(\frac1\flat\int_{\vartheta}^{\vartheta+\flat}\|\varphi(s)\|^p\,ds\right)^{1/p}.
\end{align*}
A function $\varphi$ is $W^p$-continuous if $\lim_{h\to0}\|\varphi(\cdot+h)-\varphi(\cdot)\|_{W^p}=0$, and $W^p$-bounded if $\|\varphi\|_{W^p}<\infty$; the corresponding spaces are denoted $W^p_{B}(\mathbb R,X)$ and $W^p_{BC}(\mathbb R,X)$. A function $\varphi\in W^p_{B}(\mathbb R,X)$ is $p$-th Weyl almost periodic if, for every $\epsilon>0$, the set $E\{\epsilon,\varphi\}=\{\varsigma\in\mathbb R:\|\varphi(\cdot+\varsigma)-\varphi(\cdot)\|_{W^p}<\epsilon\}$ is relatively dense in $\mathbb R$; the family of such functions is denoted by $W^p_{ap}(\mathbb R,X)$, and $W^p_{AP}(\mathbb R,X)=W^p_{ap}(\mathbb R,X)\cap W^p_{BC}(\mathbb R,X)$. The space $(W^p_{AP}(\mathbb R,X),\|\cdot\|_{W^p})$ is a complete linear space; this fact, together with translation invariance and the preservation of Weyl almost periodicity under composition with Lipschitz maps, is established in \cite[Section 2]{LiWeyl} for a general Banach space $X$, and we recall these results for reference. In Section 2.2 below we establish the corresponding statements for the square-mean spaces used throughout the remainder of the paper.

\subsection{Double-measure pseudo almost periodicity}

Let $\mathcal F$ denote the Lebesgue $\sigma$-field of $\mathbb R$ and let $\mathcal M^+$ be the set of all positive measures $\mu$ on $\mathcal F$ satisfying $\mu(\mathbb R)=+\infty$ and $\mu([a,b])<+\infty$ for all $a\le b$. Two measures $\mu,\nu\in\mathcal M^+$ are equivalent, written $\mu\sim\nu$, if there exist positive constants $\alpha_0,\beta_0$ and a bounded interval $A$ such that $\alpha_0\nu(B)\le\mu(B)\le\beta_0\nu(B)$ for all $B\in\mathcal F$ with $B\cap A=\emptyset$. For $\mu\in\mathcal M^+$ and $\sigma_0\in\mathbb R$, the shifted measure $\mu_{\sigma_0}$ is defined by $\mu_{\sigma_0}(A)=\mu(\{a+\sigma_0:a\in A\})$. We impose the following two hypotheses, standard in this theory \cite{BlotCieutatEzzinbi,DiaganaEzzinbiMiraoui}:

$(F1)$ $\mu,\nu\in\mathcal M^+$ satisfy $\limsup_{h\to\infty}\mu([-h,h])/\nu([-h,h])<\infty$.

$(F2)$ For every $\mu\in\mathcal M^+$ and $\sigma_0\in\mathbb R$, there exist a constant $\eta_0>0$ and a bounded interval $\Gamma\subset\mathbb R$ such that $\mu(a+\sigma_0:a\in A)\le\eta_0\mu(A)$ whenever $A\in\mathcal F$ and $A\cap\Gamma=\emptyset$.

Under $(F1)$ and $(F2)$, $\mu\sim\mu_{\sigma_0}$ for all $\sigma_0\in\mathbb R$; this fact, purely measure-theoretic and independent of the space $X$, is established in \cite{BlotCieutatEzzinbi} and is used below without modification. For $\mu,\nu\in\mathcal M^+$, a function $\varphi\in W^p_{BC}(\mathbb R,X)$ belongs to the ergodic class $E^p(\mathbb R,X,\mu,\nu)$ if
\begin{align*}
\lim_{h\to+\infty}\frac{1}{\nu([-h,h])}\int_{-h}^{h}\left(\lim_{\flat\to+\infty}\frac1\flat\int_{\vartheta}^{\vartheta+\flat}\|\varphi(s)\|^p\,ds\right)^{1/p}d\mu(\vartheta)=0,
\end{align*}
and a function $\varphi\in W^p_{BC}(\mathbb R,X)$ is $p$-th Weyl $(\mu,\nu)$-pseudo almost periodic if it can be written as $\varphi=\varphi_1+\varphi_2$ with $\varphi_1\in W^p_{AP}(\mathbb R,X)$ and $\varphi_2\in E^p(\mathbb R,X,\mu,\nu)$; the resulting space is denoted $W^p_{PAP}(\mathbb R,X,\mu,\nu)$. This decomposition is in general not unique \cite{LiWeyl}, a fact which is illustrated stochastically in Section 4 below.

\subsection{Square-mean Weyl almost periodicity}

For $\varphi\in L^2_{loc}(\mathbb R,H)$, the square-mean Weyl seminorm of $\varphi$ is defined by
\begin{align*}
\|\varphi\|_{W^2}=\lim_{\flat\to\infty}\sup_{\vartheta\in\mathbb R}\left(\frac1\flat\int_{\vartheta}^{\vartheta+\flat}E\|\varphi(s)\|^2\,ds\right)^{1/2}.
\end{align*}

\begin{definition}
A function $\varphi\in L^2_{loc}(\mathbb R,H)$ is square-mean $W^2$-continuous if $\lim_{h\to0}\|\varphi(\cdot+h)-\varphi(\cdot)\|_{W^2}=0$, and $W^2$-bounded if $\|\varphi\|_{W^2}<\infty$. We denote by $W^2_{B}(\mathbb R,X)$ the set of all $W^2$-bounded functions and by $W^2_{BC}(\mathbb R,X)$ the set of all $W^2$-bounded and $W^2$-continuous functions.
\end{definition}

\begin{definition}
A function $\varphi\in W^2_{B}(\mathbb R,X)$ is square-mean Weyl almost periodic if, for every $\epsilon>0$, the set $E\{\epsilon,\varphi\}=\{\varsigma\in\mathbb R:\|\varphi(\cdot+\varsigma)-\varphi(\cdot)\|_{W^2}<\epsilon\}$ is relatively dense in $\mathbb R$. The family of such functions is denoted $W^2_{ap}(\mathbb R,X)$, and we set $W^2_{AP}(\mathbb R,X)=W^2_{ap}(\mathbb R,X)\cap W^2_{BC}(\mathbb R,X)$.
\end{definition}

\begin{proposition}[Reduction principle]\label{reduction}
A function $\varphi\in L^2_{loc}(\mathbb R,X)$, regarded pointwise as an $X$-valued random variable, belongs to $L^2_{loc}(\mathbb R,H)$ if and only if $E\|\varphi(s)\|^2<\infty$ for almost every $s$ and $s\mapsto\varphi(s)$, as a map into $H$, is locally square-integrable, and in that case
\begin{align*}
\|\varphi\|_{W^2}=\|\varphi\|_{W^2_H},
\end{align*}
where the right-hand side is the deterministic Weyl seminorm of Section 2.1, taken with $p=2$ and with $X$ replaced by $H$. Consequently, $\varphi$ is square-mean $W^2$-continuous, $W^2$-bounded, or square-mean Weyl almost periodic, respectively, if and only if it is $W^2$-continuous, $W^2$-bounded, or $2$nd-Weyl almost periodic, respectively, as an $H$-valued function in the sense of Section 2.1.
\end{proposition}

\begin{proof}
Both sides of $E\|\varphi(s)\|_X^2=\|\varphi(s)\|_H^2$ are, by the definition of the norm on $H=L^2(\Omega,X)$, literally the same quantity; substituting this identity into the definitions of $\|\cdot\|_{W^2}$ (square-mean) and $\|\cdot\|_{W^2_H}$ (deterministic, $p=2$, target space $H$) shows the two seminorms coincide term by term, hence in the limit. The stated equivalences of the continuity, boundedness, and almost periodicity properties are then immediate from Definitions 2.1--2.2 above and their $H$-valued counterparts in Section 2.1.
\end{proof}

Proposition \ref{reduction} identifies the square-mean Weyl function spaces built over $X$ with the deterministic Weyl function spaces of Section 2.1 built over the Hilbert space $H=L^2(\Omega,X)$, with exponent $p=2$. Every structural property of the deterministic spaces established in \cite[Section 2]{LiWeyl} is proved there for a general Banach space, hence applies verbatim to $H$; we record the resulting statements for $W^2_{AP}(\mathbb R,X)$ below, together with self-contained proofs, both for completeness and because the composition-type statement (Theorem \ref{Lipcomp}) requires an argument beyond a formal restatement.

\begin{lemma}\label{BClinear}
If $\varphi_1,\varphi_2\in W^2_{BC}(\mathbb R,X)$ and $\alpha$ is a scalar, then $\varphi_1+\varphi_2,\ \alpha\varphi_1\in W^2_{BC}(\mathbb R,X)$.
\end{lemma}

\begin{proof}
By Minkowski's inequality for the $L^2(\Omega)$-norm, for every $\vartheta,\flat$,
\begin{align*}
\left(\frac1\flat\int_\vartheta^{\vartheta+\flat}E\|\varphi_1(s)+\varphi_2(s)\|^2\,ds\right)^{1/2}\le\left(\frac1\flat\int_\vartheta^{\vartheta+\flat}E\|\varphi_1(s)\|^2\,ds\right)^{1/2}+\left(\frac1\flat\int_\vartheta^{\vartheta+\flat}E\|\varphi_2(s)\|^2\,ds\right)^{1/2},
\end{align*}
and taking $\sup_\vartheta$ and $\flat\to\infty$ on both sides gives $\|\varphi_1+\varphi_2\|_{W^2}\le\|\varphi_1\|_{W^2}+\|\varphi_2\|_{W^2}<\infty$, so $\varphi_1+\varphi_2\in W^2_B(\mathbb R,X)$; likewise $\|\alpha\varphi_1\|_{W^2}=|\alpha|\,\|\varphi_1\|_{W^2}<\infty$. For continuity, the same triangle inequality applied to $(\varphi_1+\varphi_2)(\cdot+h)-(\varphi_1+\varphi_2)(\cdot)=[\varphi_1(\cdot+h)-\varphi_1(\cdot)]+[\varphi_2(\cdot+h)-\varphi_2(\cdot)]$ gives $\|(\varphi_1+\varphi_2)(\cdot+h)-(\varphi_1+\varphi_2)(\cdot)\|_{W^2}\le\|\varphi_1(\cdot+h)-\varphi_1(\cdot)\|_{W^2}+\|\varphi_2(\cdot+h)-\varphi_2(\cdot)\|_{W^2}\to0$ as $h\to0$; similarly for $\alpha\varphi_1$.
\end{proof}

\begin{theorem}\label{BCcomplete}
$(W^2_{BC}(\mathbb R,X),\|\cdot\|_{W^2})$ is a complete space.
\end{theorem}

\begin{proof}
Let $\{\varphi_n\}\subset W^2_{BC}(\mathbb R,X)$ be a Cauchy sequence. For every $\epsilon>0$ there is $N(\epsilon)$ such that for all $n,m\ge N$, $\|\varphi_n-\varphi_m\|_{W^2}<\epsilon$, and hence there is $\flat_0(\epsilon)>0$ such that for $\flat\ge\flat_0$,
\begin{align*}
\sup_{\vartheta\in\mathbb R}\left(\frac1\flat\int_\vartheta^{\vartheta+\flat}E\|\varphi_n(s)-\varphi_m(s)\|^2\,ds\right)^{1/2}<2\epsilon.
\end{align*}
Fixing $\flat\ge\flat_0$, this gives $\int_\vartheta^{\vartheta+\flat}E\|\varphi_n(s)-\varphi_m(s)\|^2\,ds<\flat(2\epsilon)^2$ for all $\vartheta\in\mathbb R$ and $n,m\ge N$, so $\{\varphi_n\}$ is Cauchy in $L^2([\vartheta,\vartheta+\flat],H)$ for every $\vartheta$; since $H$ is a Hilbert space, $L^2([\vartheta,\vartheta+\flat],H)$ is complete, so there is $\varphi\in L^2([\vartheta,\vartheta+\flat],H)$ with $\varphi_n\to\varphi$ there. As the limit is unique almost everywhere on any such interval, patching over $\vartheta\in\mathbb R$ gives $\varphi\in L^2_{loc}(\mathbb R,H)=L^2_{loc}(\mathbb R,X)$ (identifying random variables and elements of $H$ as in Proposition \ref{reduction}) with $\varphi_n(t)\to\varphi(t)$ in $L^2_{loc}$.

For each $\flat\ge\flat_0$ and $n\ge N$, letting $m\to\infty$ in the previous display and using Fatou's lemma,
\begin{align*}
\sup_{\vartheta\in\mathbb R}\left(\frac1\flat\int_\vartheta^{\vartheta+\flat}E\|\varphi_n(s)-\varphi(s)\|^2\,ds\right)^{1/2}\le2\epsilon,
\end{align*}
so, letting $\flat\to\infty$, $\|\varphi_n-\varphi\|_{W^2}\le2\epsilon$ for $n\ge N$. Since $\varphi-\varphi_n,\varphi_n\in W^2_B(\mathbb R,X)$ and $\varphi=(\varphi-\varphi_n)+\varphi_n$, Lemma \ref{BClinear} gives $\varphi\in W^2_B(\mathbb R,X)$. Finally, for any $h\in\mathbb R$,
\begin{align*}
\|\varphi(\cdot+h)-\varphi(\cdot)\|_{W^2}\le\|\varphi(\cdot+h)-\varphi_n(\cdot+h)\|_{W^2}+\|\varphi_n(\cdot+h)-\varphi_n(\cdot)\|_{W^2}+\|\varphi_n(\cdot)-\varphi(\cdot)\|_{W^2},
\end{align*}
and letting $n\to\infty$ then $h\to0$, using the $W^2$-continuity of each $\varphi_n$, gives $\lim_{h\to0}\|\varphi(\cdot+h)-\varphi(\cdot)\|_{W^2}=0$, so $\varphi\in W^2_{BC}(\mathbb R,X)$ and $\varphi_n\to\varphi$ in $\|\cdot\|_{W^2}$.
\end{proof}

\begin{lemma}[cf. Radov\'a \cite{Radova}]\label{commonperiod}
If $\varphi\in W^2_{ap}(\mathbb R,X)$ and $\psi\in W^2_{BC}(\mathbb R,X)\cap W^2_{ap}(\mathbb R,X)$, then $E\{\epsilon,\varphi\}\cap E\{\epsilon,\psi\}$ is relatively dense for every $\epsilon>0$.
\end{lemma}

This is the classical common-almost-period property for two almost periodic objects, valid in any complete metric space equipped with a translation action satisfying the Bochner double-sequence criterion; since the square-mean Weyl seminorm on $H$-valued functions satisfies exactly the axioms under which this property is established for a general Banach-space-valued Weyl almost periodic function in \cite{LiWeyl}, Proposition \ref{reduction} transfers the statement unchanged, and we do not reproduce its (purely metric, non-stochastic) proof here.

\begin{theorem}\label{APlinear}
If $f,g\in W^2_{AP}(\mathbb R,X)$ and $\alpha$ is a scalar, then $f+g,\ \alpha f\in W^2_{AP}(\mathbb R,X)$.
\end{theorem}

\begin{proof}
By Lemma \ref{BClinear}, $f+g,\alpha f\in W^2_{BC}(\mathbb R,X)$. Given $\epsilon>0$, by Lemma \ref{commonperiod} there is $\ell>0$ such that every interval of length $\ell$ contains $\tau$ with $\|f(\cdot+\tau)-f(\cdot)\|_{W^2}<\epsilon/2$ and $\|g(\cdot+\tau)-g(\cdot)\|_{W^2}<\epsilon/2$ simultaneously; then
\begin{align*}
\|[f(\cdot+\tau)+g(\cdot+\tau)]-[f(\cdot)+g(\cdot)]\|_{W^2}\le\|f(\cdot+\tau)-f(\cdot)\|_{W^2}+\|g(\cdot+\tau)-g(\cdot)\|_{W^2}<\epsilon,
\end{align*}
so $E\{\epsilon,f+g\}$ contains such $\tau$'s and is relatively dense. That $\alpha f\in W^2_{ap}(\mathbb R,X)$ follows from $E\{\epsilon/|\alpha|,f\}\subset E\{\epsilon,\alpha f\}$ for $\alpha\ne0$ (trivial for $\alpha=0$).
\end{proof}

\begin{theorem}\label{APtrans}
If $x\in W^2_{AP}(\mathbb R,X)$ and $\sigma\in\mathbb R$, then $x(\cdot+\sigma)\in W^2_{AP}(\mathbb R,X)$; that is, $W^2_{AP}(\mathbb R,X)$ is translation invariant.
\end{theorem}

\begin{proof}
For any $\sigma,\tau\in\mathbb R$,
\begin{align*}
\|x(\cdot+\sigma+\tau)-x(\cdot+\sigma)\|_{W^2}
&=\lim_{\flat\to\infty}\sup_{\vartheta\in\mathbb R}\left(\frac1\flat\int_\vartheta^{\vartheta+\flat}E\|x(t+\sigma+\tau)-x(t+\sigma)\|^2\,dt\right)^{1/2}\\
&=\lim_{\flat\to\infty}\sup_{\vartheta\in\mathbb R}\left(\frac1\flat\int_{\vartheta+\sigma}^{\vartheta+\sigma+\flat}E\|x(t+\tau)-x(t)\|^2\,dt\right)^{1/2}\\
&\le\lim_{\flat\to\infty}\sup_{\vartheta'\in\mathbb R}\left(\frac1\flat\int_{\vartheta'}^{\vartheta'+\flat}E\|x(t+\tau)-x(t)\|^2\,dt\right)^{1/2}=\|x(\cdot+\tau)-x(\cdot)\|_{W^2},
\end{align*}
using the substitution $t\mapsto t+\sigma$ and the fact that $\vartheta+\sigma$ ranges over all of $\mathbb R$ as $\vartheta$ does. Hence $E\{\epsilon,x\}\subset E\{\epsilon,x(\cdot+\sigma)\}$, and $x(\cdot+\sigma)\in W^2_{ap}(\mathbb R,X)$; $W^2$-boundedness and continuity of $x(\cdot+\sigma)$ are immediate from those of $x$.
\end{proof}

\begin{theorem}\label{Lipcomp}
If $g:X\to X$ satisfies $\|g(x)-g(y)\|\le L_g\|x-y\|$ for all $x,y\in X$, and $\varphi\in W^2_{AP}(\mathbb R,X)$, then $t\mapsto g(\varphi(t))$, defined pathwise by $g(\varphi(t))(\omega):=g(\varphi(t,\omega))$, belongs to $W^2_{AP}(\mathbb R,X)$.
\end{theorem}

\begin{proof}
For each $t$, $g\circ\varphi(t)$ is a well-defined element of $H$ since $E\|g(\varphi(t))\|^2\le2L_g^2E\|\varphi(t)\|^2+2\|g(0)\|^2<\infty$. Regarding $g$ as inducing the Nemytskii map $\hat g:H\to H$, $\hat g(\xi)(\omega):=g(\xi(\omega))$, we have
\begin{align*}
\|\hat g(\xi)-\hat g(\eta)\|_H^2=E\|g(\xi)-g(\eta)\|^2\le L_g^2E\|\xi-\eta\|^2=L_g^2\|\xi-\eta\|_H^2,
\end{align*}
so $\hat g$ is Lipschitz on $H$ with the same constant $L_g$. By Minkowski's inequality, for every $\vartheta,\flat$,
\begin{align*}
\left(\frac1\flat\int_\vartheta^{\vartheta+\flat}\|\hat g(\varphi(t))\|_H^2\,dt\right)^{1/2}\le L_g\left(\frac1\flat\int_\vartheta^{\vartheta+\flat}\|\varphi(t)\|_H^2\,dt\right)^{1/2}+\|g(0)\|,
\end{align*}
so $\|\hat g\circ\varphi\|_{W^2}\le L_g\|\varphi\|_{W^2}+\|g(0)\|<\infty$, and, for any $\tau\in\mathbb R$,
\begin{align*}
\|\hat g(\varphi(\cdot+\tau))-\hat g(\varphi(\cdot))\|_{W^2}\le L_g\|\varphi(\cdot+\tau)-\varphi(\cdot)\|_{W^2},
\end{align*}
so $E\{\epsilon/L_g,\varphi\}\subset E\{\epsilon,\hat g\circ\varphi\}$ (with the convention $E\{\epsilon/L_g,\varphi\}=\mathbb R$ if $L_g=0$), whence $\hat g\circ\varphi\in W^2_{ap}(\mathbb R,X)$; the same estimate with $\tau=h\to0$ gives $W^2$-continuity, so $\hat g\circ\varphi\in W^2_{AP}(\mathbb R,X)$.
\end{proof}

\begin{theorem}\label{APconv}
If $\{\varphi_n\}\subset W^2_{AP}(\mathbb R,X)$ and $\varphi_n\to\varphi$ with respect to $\|\cdot\|_{W^2}$, then $\varphi\in W^2_{AP}(\mathbb R,X)$.
\end{theorem}

\begin{proof}
By Theorem \ref{BCcomplete}, $\varphi\in W^2_{BC}(\mathbb R,X)$, and for any $\epsilon>0$ there is $N$ with $\|\varphi_N-\varphi\|_{W^2}<\epsilon$. Since $\varphi_N\in W^2_{AP}(\mathbb R,X)$, there is $\ell(\epsilon)>0$ such that every interval of length $\ell$ contains $\tau$ with $\|\varphi_N(\cdot+\tau)-\varphi_N(\cdot)\|_{W^2}<\epsilon$. For such $\tau$,
\begin{align*}
\|\varphi(\cdot+\tau)-\varphi(\cdot)\|_{W^2}\le\|\varphi(\cdot+\tau)-\varphi_N(\cdot+\tau)\|_{W^2}+\|\varphi_N(\cdot+\tau)-\varphi_N(\cdot)\|_{W^2}+\|\varphi_N(\cdot)-\varphi(\cdot)\|_{W^2}<3\epsilon,
\end{align*}
using Theorem \ref{APtrans} to note $\|\varphi(\cdot+\tau)-\varphi_N(\cdot+\tau)\|_{W^2}=\|(\varphi-\varphi_N)(\cdot+\tau)\|_{W^2}\le\|\varphi-\varphi_N\|_{W^2}<\epsilon$ (translation invariance of the seminorm itself, immediate from the definition of $\sup_\vartheta$). Hence $E\{3\epsilon,\varphi\}\supset E\{\epsilon,\varphi_N\}$ is relatively dense, and $\varphi\in W^2_{ap}(\mathbb R,X)$.
\end{proof}

\begin{theorem}\label{APcomplete}
$(W^2_{AP}(\mathbb R,X),\|\cdot\|_{W^2})$ is a complete linear space.
\end{theorem}

\begin{proof}
By Theorem \ref{APlinear}, $W^2_{AP}(\mathbb R,X)$ is a linear space. By Theorem \ref{BCcomplete}, $W^2_{BC}(\mathbb R,X)$ is complete, and by Theorem \ref{APconv}, $W^2_{AP}(\mathbb R,X)$ is a closed subset of $W^2_{BC}(\mathbb R,X)$; a closed subset of a complete metric space is complete.
\end{proof}

Given $\mu,\nu\in\mathcal M^+$ satisfying $(F1)$ and $(F2)$, we define the square-mean ergodic class
\begin{align*}
E^2(\mathbb R,X,\mu,\nu):\qquad \lim_{h\to+\infty}\frac{1}{\nu([-h,h])}\int_{-h}^{h}\left(\lim_{\flat\to+\infty}\frac1\flat\int_{\vartheta}^{\vartheta+\flat}E\|\varphi(s)\|^2\,ds\right)^{1/2}d\mu(\vartheta)=0,
\end{align*}
and the space of square-mean Weyl $(\mu,\nu)$-pseudo almost periodic functions $W^2_{PAP}(\mathbb R,X,\mu,\nu)=W^2_{AP}(\mathbb R,X)\oplus E^2(\mathbb R,X,\mu,\nu)$. We now establish, in detail, the three structural properties of $E^2(\mathbb R,X,\mu,\nu)$ used later: an equivalent characterization, invariance under equivalent measures, and translation invariance.


\begin{theorem}\label{E2char}
Let $\mu,\nu\in\mathcal M^+$, let $\Gamma$ be a bounded interval (possibly empty), and assume $(F1)$ and $\varphi\in W^2_{BC}(\mathbb R,X)$. Then the following assertions are equivalent:
\begin{align*}
&\text{(i)}\quad \varphi\in E^2(\mathbb R,X,\mu,\nu);\\
&\text{(ii)}\quad \lim_{h\to+\infty}\frac1{\nu([-h,h]\setminus\Gamma)}\int_{[-h,h]\setminus\Gamma}\left(\lim_{\flat\to+\infty}\frac1\flat\int_\vartheta^{\vartheta+\flat}E\|\varphi(s)\|^2\,ds\right)^{1/2}d\mu(\vartheta)=0;\\
&\text{(iii)}\quad \text{for every }\epsilon>0,\ \lim_{h\to\infty}\frac{\mu(\{\vartheta\in[-h,h]\setminus\Gamma:(E_\flat\|\varphi\|^2(\vartheta))^{1/2}>\epsilon\})}{\nu([-h,h]\setminus\Gamma)}=0,
\end{align*}
where $(E_\flat\|\varphi\|^2(\vartheta))^{1/2}:=\left(\lim_{\flat\to+\infty}\frac1\flat\int_\vartheta^{\vartheta+\flat}E\|\varphi(s)\|^2\,ds\right)^{1/2}$.
\end{theorem}

\begin{proof}
Write $A=\nu(\Gamma)$, $C=\mu(\Gamma)$, and $B=\int_\Gamma(E_\flat\|\varphi\|^2(\vartheta))^{1/2}\,d\mu(\vartheta)<\infty$ (finite since $\Gamma$ is bounded and $\varphi\in W^2_{BC}(\mathbb R,X)$).

(i)$\Rightarrow$(ii). Since $\int_{[-h,h]\setminus\Gamma}=\int_{[-h,h]}-\int_\Gamma$ and $\nu([-h,h]\setminus\Gamma)=\nu([-h,h])-A$,
\begin{align*}
\frac1{\nu([-h,h]\setminus\Gamma)}\int_{[-h,h]\setminus\Gamma}(E_\flat\|\varphi\|^2)^{1/2}d\mu
=\frac{\nu([-h,h])}{\nu([-h,h])-A}\left(\frac1{\nu([-h,h])}\int_{[-h,h]}(E_\flat\|\varphi\|^2)^{1/2}d\mu-\frac{B}{\nu([-h,h])}\right).
\end{align*}
Since $\nu(\mathbb R)=\infty$, $\nu([-h,h])\to\infty$, so $\nu([-h,h])/(\nu([-h,h])-A)\to1$ and $B/\nu([-h,h])\to0$; hence (ii) is equivalent to $\lim_{h\to\infty}\frac1{\nu([-h,h])}\int_{[-h,h]}(E_\flat\|\varphi\|^2)^{1/2}d\mu=0$, which is (i).

(iii)$\Rightarrow$(ii). Set $A_{\epsilon,h}=\{\vartheta\in[-h,h]\setminus\Gamma:(E_\flat\|\varphi\|^2(\vartheta))^{1/2}>\epsilon\}$ and $B_{\epsilon,h}=([-h,h]\setminus\Gamma)\setminus A_{\epsilon,h}$. Splitting the integral in (ii) over $A_{\epsilon,h}$ and $B_{\epsilon,h}$, we get
\begin{align*}
\frac1{\nu([-h,h]\setminus\Gamma)}\int_{[-h,h]\setminus\Gamma}(E_\flat\|\varphi\|^2)^{1/2}d\mu
\le2\|\varphi\|_{W^2}\frac{\mu(A_{\epsilon,h})}{\nu([-h,h]\setminus\Gamma)}+\epsilon\,\frac{\mu([-h,h]\setminus\Gamma)}{\nu([-h,h]\setminus\Gamma)}.
\end{align*}
Writing $\mu([-h,h]\setminus\Gamma)/\nu([-h,h]\setminus\Gamma)=\left(\mu([-h,h])/\nu([-h,h])\right)\cdot\left(1-C/\mu([-h,h])\right)/\left(1-A/\nu([-h,h])\right)$, and using $\mu(\mathbb R)=\nu(\mathbb R)=\infty$ together with $(F1)$, the right factor tends to $1$ and \\ $\limsup_h\mu([-h,h])/\nu([-h,h])=:K_0<\infty$; hence, if (iii) holds, letting $h\to\infty$ and then $\epsilon\to0$,
\begin{align*}
\limsup_{h\to\infty}\frac1{\nu([-h,h]\setminus\Gamma)}\int_{[-h,h]\setminus\Gamma}(E_\flat\|\varphi\|^2)^{1/2}d\mu\le K_0\epsilon,
\end{align*}
for every $\epsilon>0$, giving (ii).

(ii)$\Rightarrow$(iii). From
\begin{align*}
\frac1{\nu([-h,h]\setminus\Gamma)}\int_{[-h,h]\setminus\Gamma}(E_\flat\|\varphi\|^2)^{1/2}d\mu\ge\frac1{\nu([-h,h]\setminus\Gamma)}\int_{A_{\epsilon,h}}(E_\flat\|\varphi\|^2)^{1/2}d\mu\ge\epsilon\,\frac{\mu(A_{\epsilon,h})}{\nu([-h,h]\setminus\Gamma)},
\end{align*}
(ii) forces $\mu(A_{\epsilon,h})/\nu([-h,h]\setminus\Gamma)\to0$ as $h\to\infty$, for every $\epsilon>0$, which is (iii).
\end{proof}

\begin{thebibliography}{99}
\bibitem[BlotCieutatEzzinbi]{BlotCieutatEzzinbi} J. Blot, P. Cieutat, K. Ezzinbi, \emph{Measure theory and pseudo almost automorphic functions: New developments and applications}, Nonlinear Anal. \textbf{75}, (2012), 2426--2447.
\bibitem[DiaganaEzzinbiMiraoui]{DiaganaEzzinbiMiraoui} T. Diagana, K. Ezzinbi, M. Miraoui, \emph{Pseudo-almost periodic and pseudo-almost automorphic solutions to some evolution equations involving theoretical measure theory}, Cubo \textbf{16}(2), (2014), 1--32.
\bibitem[LiWeyl]{LiWeyl} Y. Li, \emph{Weyl double-measure pseudo almost periodic functions and Weyl double-measure pseudo almost periodic solutions to semilinear evolution equations}, Quaest. Math. \textbf{49}(6), (2026), 777--799.
\bibitem[Radova]{Radova} L. Radov\'a, \emph{Theorems of Bohr-Neugebauer-type for almost-periodic differential equations}, Math. Slovaca \textbf{54}(2), (2004), 191--207.
\end{thebibliography}
\end{document}
